% ============================================================
%  Poul Martin Møller,
%  "Hvad er Philosophie?"
%
%  SOURCE: Introduction to Møller's lecture course on the history
%  of older philosophy, delivered at the University of Copenhagen
%  in 1834--35. The text is a composite: Møller's own manuscript
%  (written "with a brisk pen") supplemented from students' lecture
%  notes by editor F.C. Olsen. The editorial notes in the text are
%  student additions, not Møller's own.
%
%  The introductory section ("Indledning") is singled out by Børge
%  Madsen (ed. of Filosofiske Essays og Strøtanker) for its "freer
%  form, freshness, and many striking philosophical observations,"
%  distinguishing it from the heavier historical material that
%  follows.
%
%  Søren Kierkegaard was almost certainly a student in this course.
%  He enrolled at the University of Copenhagen in 1830; Møller
%  became professor there in 1831; Kierkegaard is documented as
%  having attended Møller's lectures and held him as his most
%  important teacher.
%
%  Published posthumously in:
%  Efterladte Skrifter af Poul M. Møller, bd. 3,
%  ed. F.C. Olsen. Copenhagen: C.A. Reitzel, 1843.
%
%  Also in: Filosofiske Essays og Strøtanker,
%  udvalg og indledning ved Børge Madsen.
%  [Publisher and date to be confirmed from title page.]
%
%  TRANSCRIPTION (Danish)
%  Completed: Indledning (pp. 115--118 in Børge Madsen edition)
% ============================================================

\documentclass[12pt]{article}
\usepackage[utf8]{inputenc}
\usepackage[T1]{fontenc}
\usepackage{libertinus}
\usepackage{libertinust1math}
\usepackage[danish]{babel}
\usepackage[protrusion=true,expansion=true]{microtype}
\usepackage[margin=1.4in]{geometry}
\usepackage{setspace}
\usepackage{fancyhdr}
\usepackage[colorlinks=true,linkcolor=black,urlcolor=black,
  pdftitle={Hvad er Philosophie?},
  pdfauthor={Poul Martin Møller}]{hyperref}
\setlength{\headheight}{15pt}
\pagestyle{fancy}
\fancyhf{}
\fancyhead[LE,RO]{\thepage}
\fancyhead[RE]{\textit{Hvad er Philosophie?}}
\fancyhead[LO]{\textit{\rightmark}}
\setlength{\parindent}{1.5em}
\setlength{\parskip}{0pt}
\onehalfspacing

\title{\textbf{Hvad er Philosophie?}\\[0.5em]
  \large Indledning til Forelæsninger over den ældre Filosofis Historie\\
  Universitetet i Kjøbenhavn, 1834--35}
\author{Poul Martin Møller}
\date{Udgivet posthumt i \textit{Efterladte Skrifter}, bd.~3,\\
  red.\ F.C.\ Olsen. Kjøbenhavn: C.A.\ Reitzel, 1843}

\begin{document}
\maketitle

% ============================================================
%  INDLEDNING
% ============================================================

% pp. 115--118 (Børge Madsen edition; page range in
% Efterladte Skrifter bd. 3 to be confirmed)

\noindent
Hvad er Philosophie? Dette Spørgsmaal kunde man vente besvaret
af den, der vil foredrage Philosophiens Historie. Det er let at
give en heel Deel lige rigtige Definitioner paa Philosophie, men
de ville kun tjene til at minde den, der veed, hvad Philosophie
er, om, hvad han tænker sig derved. Dersom han ikke i Forveien
havde et Begreb eller en Forestilling om Philosophien, vilde han
ikke ved at høre en enkelt Sætning med Eet komme til at
erkjende, hvad Philosophie er. Man kom nemlig til at bestemme
Philosophien ved Hjelp af Begreber, som stode i nær Forbindelse
dermed; men dem kan Ingen have paa en grundig Maade, uden at han
har philosopheret og altsaa veed, hvad Philosophie er. Jeg
sætter, at jeg vilde sige: Philosophien er en Videnskab om
Tingenes sidste Grunde eller om det Evige, da vilde dog Ingen,
som slet ikke vidste, hvad Philosophie er, kunne tydelig vide,
hvad jeg mener ved Tingenes sidste Grunde eller det Evige.

Hertil kommer endnu en Vanskelighed. Enhver Sætning, i hvilken
man vilde udtrykke Philosophiens Begreb, er saa eensidig, at man
ved at tænke Noget derved strax føres til at forandre sit Udtryk
og at fremføre modsatte Bestemmelser. Jeg siger, Philosophien er
Videnskaben om det Evige; men da den ikke skal afhandle det
abstracte Evige, men det sig aabenbarende Evige, kan jeg med
ligesaa megen Sandhed sige, at den er Videnskaben om det
Endelige. --- Jeg siger, at den er Videnskaben om Tingenes
sidste Grunde; men den er ligesaavelVidenskaben om disse
Grundes Følge, da Grunde først i deres Følge vise sig som
Grunde. Saaledes drives jeg ved ethvert eensidigt Prædicat
videre frem og kan kun med Vilkaarlighed standse ved en
eensidig Definition. Det fuldstændige Udtryk for Philosophiens
Begreb er Fremstillingen af Philosophien selv, og enhver
Forklaring, der standser, inden denne hele Fremstilling er
fuldbragt, er egentlig kun tildeels rigtig. Der er saaledes
ikke vundet Noget ved den sædvanlige, formeentlig nøiagtige
og videnskabelige Maade at forudskikke en Definition. Denne
Formalitet overspringes altsaa, og der forudsættes, hvad der
naturligviis med Grund kan forudsættes --- at Den, der hører
disse Forelæsninger, veed hvad Philosophie er, ligesaa godt
som efterat have hørt hvilkensomhelst Definition.

Derimod kunde det maaskee være passende at tale Noget om
Forskjellen imellem Philosophien paa den ene Side, og paa
den anden Side Religionslæren og den videnskabelige Tænkning,
der ikke er philosophisk. Men det tør jeg kun gjøre med et
Par Ord, da de Distinctioner, som her ere at gjøre, i den
senere Tid saa ofte have været gjorte, at de snart maae
forekomme Enhver, der blot har lest endeel tydske og danske
Recensioner og Fortaler til philosophiske Afhandlinger, som
gamle velbekjendte Ramser.

For det Første er Philosophien forskjellig fra en Religionslære,
der ikke er philosophisk. De ere ikke forskjellige i deres
Indhold, da Gud og hans Forhold til Verden ere begges fælles
Hovedgjenstande. Men Religionslæren kan i sin Form være
forskjellig fra Philosophien, og den kan være det meer eller
mindre. Jo mere den religiøse Bevidstheds Indhold blot er
fremstillet paa en symbolsk Maade, jo mere den istedetfor rene
Tankebestemmelser sætter Phantasiebilleder og betydningsfulde
Myther, desto meer er den forskjellig fra Philosophien. Det
fornuftige Væsen, der med Tro modtager billedlige Fremstillinger
af Guddommens Væsen og sindrige Fortællinger om Gudernes Liv og
Verdens Oprindelse, modtager Noget meer eller mindre fremmed,
Noget, som han ikke kan finde paa samme Maade hos sig selv,
fordi det, der bydes ham, ikke er almeengyldige Frembringelser
af den fornuftige Tænkning, der hos Alle er eens. Sandselige
Forestillinger om det Evige høre ikke med til Philosophien,
om der end ligger en nok saa dybsindig Verdensanskuelse til
Grund derfor. Dette bemærkes fordi man ofte har villet drage
saadanne Forestillingers Historie ind under Philosophiens
Historie, især hvor Talen har været om Orientens Philosophie.
Philosophien har kun med Tanker at gjøre; om end de
dybsindigste Tanker ligge skjulte i Myther og Phantasiebilleder,
befinder man sig her udenfor Philosophiens Territorium. Thi
der skulle Tankerne ikke ligge skjulte, men bestemt udvikles.
Derimod kunne sindrige Myther og religiøse Forestillinger bane
Veien for Philosophien, tjene til Incitamenter for den
philosophiske Tænkning, og det lærer Historien, at de ofte
have gjort det.

Mere nærmer Religionslæren sig til Philosophien, naar den
indeholder Tanker i Form af Tanker; men naar de religiøse
Sætninger blot antages i Tro og Tillid til dens Ufeilbarhed,
der har overleveret dem, da kan heller ikke denne Slags Lære,
i hvor sand og dyb den end forresten kan være, regnes til
philosophisk Fremstilling af Sandheden. Dertil hører, at
Tænkningen selv skal være den eneste Erkjendelseskilde, til
hvilken der provoceres. Fremstillingen af Sandheden maa være
almeengyldig for at være philosophisk, saa at ethvert tænkende
Individ, der følger Fremstillingen, kan følge den med fuld
Overbeviisning, fordi den er et Udtryk af den Tankegang, som
er nødvendig bestemt af den universelle Fornuft. Philosophien
er da forskjellig fra den Religionslære, der beraaber sig paa
en anden Autoritet end Fornuften selv, enten denne bevæger sig
paa Forestillingens eller paa Tankernes Enemærker.

Philosophien maa fremdeles adskilles fra den videnskabelige
Tænkning, som ikke er philosophisk. Denne har ikke sit Indhold
tilfælles med Philosophien, men den har det tilfælles med
Philosophien, at den gjør den selvvirksomme Tænkning gjeldende,
og betragter den som en Vei til Sandhedserkjendelse. Den har
ogsaa det tilfælles med Philosophien, at den gaaer ud over det,
der er givet ved Sandsning, og stræber at erkjende det
Almindelige, der ligger til Grund for Phænomenerne, henfører
dem under almindelige Begreber og stræber at erkjende de Love,
der gjelde for deres Forandringer. Men for saavidt som den er
en Tænkning over noget Givet og gaaer ud paa at ordne og
forbinde en i Sandseverdenen forefunden Sum af Forestillinger
om tilværende Ting, er det endnu ikke den egentlig philosophiske
Tænkning. --- Den videnskabelige Tænkning, der falder udenfor
Philosophien, viser vel hen til den philosophiske Idee, da en
dunkel Forestilling om Eenheden i den hele Tilværelse leder
alle dens Bestræbelser; men den søger at komme til
Erkjendelsen af denne Eenhed ved at sammenligne og forbinde
det Endelige under mere almindelige Overblik eller vel endog
ved at udelade den sandselige Bestemthed, at abstrahere fra
den. --- Under sin hele Virksomhed gjør den en regulativ Brug
af Tænkningens egne Begreber, de aprioriske Begreber eller
Kategorier; for saavidt er den tildeels ved sig selv eller
har med sig selv at gjøre, foruden med det udenfra givne
Stof; men den er ikke kommen til Bevidsthed om disse Begrebers
Natur og Væsen og deres Forhold til Tingenes Natur og Væsen.
Deraf bliver deels Følgen, 1) at den bruger disse Begreber
paa Slump, uden at agte paa deres Forhold til hinanden: en
Mængde Kategorier bruges efter en Erindring om, hvorledes
Andre bruge dem; men de betragtes som faste Forudsætninger,
som Forstandens Dine, der skal sees igjennem, men som ikke
selv skulle sees. 2) Dernæst bliver Følgen deraf, at man vel
mærker sin Tænknings consequente Fremgangsmaade og dens
Forudsætning af Begreber, der ikke ere givne ved Sandsning,
men antager dem for at være blot subjective og ikke gjeldende
for Tingene.

I den philosophiske Tænkning derimod kommer Tænkningen til
Bevidsthed om, at dens Natur og Væsen er Tilværelsens Natur
og Væsen, og at den ved at tænke blot efter sine egne Love
erkjender Tilværelsens evige og sande Bestemmelser. Den
kommer til Bevidsthed om, at, hvad der nødvendigviis maa
tænkes, ogsaa er saaledes, som man nødes til at tænke det.
Den Idee [Eenhed i den hele Tilværelse], der blot tjente til
regulativt Brug under den empiriske Tænkning, erkjendes da
for at være det rette Udgangspunct og det ubestemte
Begyndelsespunct, som man nærmere maa bestemme ved
Tænkningens egen Fuldmagt, for at komme til en fornuftig
Erkjendelse af Tingene.\footnote{Dette er Grundtanken i
  enhver speculativ Logik (eller forhen Metaphysik, Ontologie,
  Colleg.).}

Saaledes have vi paa en relativ Maade begrendset Philosophiens
Begreb til to Sider og udelukket adskillige Tankeproducter
fra det aandelige Indhold, som vi ville agte paa i Udviklingen
af Philosophiens Historie.

Philosophiens Historie hænger naturligviis sammen med
Menneskehedens almindelige Historie. Den frie Tænkning er
paa mange Maader bleven modificeret, hindret og fremmet, ved
de udvortes Forhold, hvorunder den har udviklet sig. Den, der
vil tilfulde forklare sig den Skikkelse, den har antaget hos
et givet Folk til en given Tid, behøver derfor ikke sjelden
Kundskaber til Folkets Charakteer og Historie. Dette ville
vi udvikle noget nærmere.

Philosophien --- saadan som den existerer og træder virkelig
frem i en menneskelig Bevidsthed --- har altid en bestemt
Sammenhæng med Philosophen, med hans hele aandelige Liv og
Forestillingskreds; men denne er igjen for det Første for en
stor Deel betinget ved den naturlige Individualitet og det
Dannelsestrin, der er ---

% [text to be continued; remainder of lecture course to be added]


% ============================================================
%  RESTEN AF FORELÆSNINGERNE
% ============================================================

% The remainder of the lecture course covers the history of
% older (ancient and medieval) philosophy. Sections to be
% identified and added from Efterladte Skrifter bd. 3.

% [text to be added]

\end{document}
